% Part: turing-machines
% Chapter: undecidability
% Section: enumerating-tms

\documentclass[../../../include/open-logic-section]{subfiles}

\begin{document}

\olfileid{tur}{und}{enu}
\olsection{ట్యూరింగ్ యంత్రాలను లెక్కించడం}

\begin{explain}
అన్ని ట్యూరింగ్ యంత్రాల సమితి \tetoken{లెక్కించదగిన}{!!{enumerable}} అని చూపవచ్చు. ప్రతి ట్యూరింగ్
యంత్రాన్నీ పరిమితంగా వర్ణించవచ్చనే వాస్తవం నుండి ఇది అనుగమిస్తుంది. స్థితుల
సమితి, టేపు వర్ణమాల పరిమిత సమితులు. సంక్రమణ ప్రమేయం $Q \times \Sigma$
నుండి $Q \times \Sigma \times \{\TMleft, \TMright, \TMstay\}$కు వెళ్లే
పాక్షిక ప్రమేయం; కాబట్టి అది నిర్వచితమైన పరిమిత సంఖ్యలోని ఆర్గ్యుమెంట్
యుగ్మాలకు దాని విలువలను జాబితా చేసి దానినీ నిర్దేశించవచ్చు.

ఇంతవరకు ఇది నిజమే, కానీ ఒక సూక్ష్మమైన తేడా ఉంది. స్థితుల సమితిలో, టేపు
వర్ణమాలలో ఎలాంటి \tetoken{మూలకాలు}{!!{element}s} ఉండవచ్చనే దానిపై ట్యూరింగ్ యంత్రాల నిర్వచనం
ఏ పరిమితినీ విధించలేదు. కాబట్టి, ఉదాహరణకు, ప్రతి వాస్తవ సంఖ్యకూ ఆ సంఖ్యను
ఒక స్థితిగా ఉపయోగించే ట్యూరింగ్ యంత్రం సాంకేతికంగా ఉంది. అయితే ట్యూరింగ్
యంత్రపు \emph{ప్రవర్తన}, ఏ వస్తువులు స్థితులుగా, వర్ణమాలగా పనిచేస్తాయనే
దానిపై ఆధారపడదు.
\olref{fig:variants}లోని రెండు ట్యూరింగ్ యంత్రాలను పరిశీలించండి.
\begin{figure}
\begin{center}
\begin{tikzpicture}[->,>=stealth',shorten >=1pt,auto,node distance=2.8cm,
                    semithick]
  \tikzstyle{every state}=[fill=none,draw=black,text=black]

  \node[initial,state]         (A)              {$q_0$};
  \node[state]         (B) [right of=A] {$q_1$};

  \path (A) edge [bend left] node {\TMtrans{\TMstroke}{\TMstroke}{\TMright}} (B)
        (B) edge [loop above] node {\TMtrans{\TMblank}{\TMblank}{\TMright}} (B)
            edge [bend left] node {\TMtrans{\TMstroke}{\TMstroke}{\TMright}} (A);
\end{tikzpicture}\\
\begin{tikzpicture}[->,>=stealth',shorten >=1pt,auto,node distance=2.8cm,
  semithick]
\tikzstyle{every state}=[fill=none,draw=black,text=black]

\node[initial,state]         (A)              {$s$};
\node[state]         (B) [right of=A] {$h$};

\path (A) edge [bend left] node {\TMtrans{A}{A}{\TMright}} (B)
(B) edge [loop above] node {\TMtrans{\TMblank}{\TMblank}{\TMright}} (B)
edge [bend left] node {\TMtrans{A}{A}{\TMright}} (A);
\end{tikzpicture}
\end{center}
\caption{\emph{సరి} యంత్రపు రూపాంతరాలు}
\ollabel{fig:variants}
\end{figure}
ఈ రెండు రేఖాచిత్రాలు రెండు యంత్రాలకు అనుగుణంగా ఉన్నాయి: టేపు వర్ణమాల
$\Sigma = \{\TMendtape,\TMblank,\TMstroke\}$, స్థితుల సమితి
$\{q_0,q_1\}$ గల $M$; వర్ణమాల $\Sigma' =
\{\TMendtape,\TMblank,A\}$, స్థితులు $\{s,h\}$ గల $M'$. కానీ వాటి
నిర్దేశాలు మిగతా విధంగా ఒకటే: $n$ సరి అయితే, అప్పుడే మాత్రమే $n$
$\TMstroke$ల క్రమంపై $M$ ఆగుతుంది; $n$ సరి అయితే, అప్పుడే మాత్రమే $n$
$A$ల క్రమంపై $M'$ ఆగుతుంది. మనం చేసినదల్లా $\TMstroke$కు~$A$ అనీ,
$q_0$కు~$s$ అనీ, $q_1$కు~$h$ అనీ పేర్లు మార్చడం. ఈ ఉదాహరణను
సాధారణీకరించవచ్చు: సంకేతాలు, స్థితులకు ఇలాంటి పేరు మార్పు వల్ల ఒక యంత్రం
నుండి మరొకటి వస్తే, ఆ ట్యూరింగ్ యంత్రాలను ఒకటిగానే భావించవచ్చు. నిజానికి,
ట్యూరింగ్ యంత్రపు సంకేతాలు, స్థితులను ధన పూర్ణ సంఖ్యలుగానే భావించవచ్చు:
$\sigma_0$కు బదులు~$1$, $\sigma_1$కు బదులు~$2$, తదితరాలు;
$\TMendtape$కు~$1$, $\TMblank$కు~$2$, తదితరాలు. ఈ విధంగా \emph{సరి}
యంత్రం
\olref{fig:standard-even}లో చూపిన యంత్రంగా మారుతుంది.
\begin{figure}
\[\begin{tikzpicture}[->,>=stealth',shorten >=1pt,auto,node distance=2.8cm,
  semithick]
\tikzstyle{every state}=[fill=none,draw=black,text=black]

\node[initial,state]         (A)              {$1$};
\node[state]         (B) [right of=A] {$2$};

\path (A) edge [bend left] node {\TMtrans{3}{3}{\TMright}} (B)
(B) edge [loop above] node {\TMtrans{2}{2}{\TMright}} (B)
edge [bend left] node {\TMtrans{3}{3}{\TMright}} (A);
\end{tikzpicture}
\]
\caption{ఒక ప్రామాణిక \emph{సరి} యంత్రం}
\ollabel{fig:standard-even}
\end{figure}
స్థితులు, సంకేతాలు ధన పూర్ణ సంఖ్యలైన ట్యూరింగ్ యంత్రాన్ని \emph{ప్రామాణిక}
యంత్రం అని పిలిచి, ఇకపై ప్రామాణిక యంత్రాలను మాత్రమే పరిశీలించవచ్చు.
\footnote{``ప్రామాణిక యంత్రం'' అనే పదప్రయోగం ప్రామాణికమైనది కాదు.}

ట్యూరింగ్ యంత్రాల సమితి \tetoken{లెక్కించదగిన}{!!{enumerable}} అని చూపాలనుకున్నాం; పై పరిశీలనలను
దృష్టిలో ఉంచుకుంటే, ప్రామాణిక ట్యూరింగ్ యంత్రాల సమితి \tetoken{లెక్కించదగిన}{!!{enumerable}} అని
చూపితే చాలు. ఒక ప్రామాణిక ట్యూరింగ్ యంత్రం
$M = \tuple{Q, \Sigma, q_0, \delta}$ ఇచ్చారని అనుకుందాం. ధన పూర్ణ సంఖ్యల
పరిమిత సంకేతమాలతో దాన్ని ఎలా వర్ణించగలం? మొదట స్థితుల సంఖ్య, స్థితులు,
సంకేతాల సంఖ్య, సంకేతాలు, ప్రారంభ స్థితిని జాబితా చేస్తాం. ($M$ ప్రామాణిక
యంత్రం కనుక ఇవన్నీ ధన పూర్ణ సంఖ్యలని గుర్తుంచుకోండి.) మరి~$\delta$ సంగతేమిటి?
సాధ్యమైన ఆర్గ్యుమెంట్ల సమితి, అంటే $\tuple{q,\sigma}$ యుగ్మాల సమితి,
$Q$,~$\Sigma$ పరిమితాలు కనుక పరిమితమే. అందువల్ల~$\delta$లోని సమాచారం,
$\delta(q,\sigma) = \tuple{q', \sigma', D}$ అయ్యే అన్ని
$5$-టపుళ్ల $\tuple{q, \sigma, q', \sigma', d}$ పరిమిత జాబితా మాత్రమే;
ఇక్కడ $d$ అనేది దిశ~$D$ను సంకేతీకరించే సంఖ్య (ఉదాహరణకు,~$\TMleft$కు~$1$,
$\TMright$కు~$2$, $\TMstay$కు~$3$).

ఈ విధంగా ప్రతి ప్రామాణిక ట్యూరింగ్ యంత్రాన్ని ధన పూర్ణ సంఖ్యల పరిమిత
జాబితాగా, అంటే $s_M \in (\PosInt)^*$ అనే క్రమంగా వర్ణించవచ్చు. ఉదాహరణకు,
ప్రామాణిక \emph{సరి} యంత్రం కింది క్రమంతో సంకేతీకరించబడుతుంది:
\[
2, \underbrace{1, 2}_Q, 3, \overbrace{1, 2, 3}^\Sigma, 1, \underbrace{1, 3, 2, 3, 2}_{\delta(1,3) = \tuple{2,3,R}},
\overbrace{2, 2, 2, 2, 2}^{\delta(2,2) = \tuple{2,2,R}},
\underbrace{2, 3, 1, 3, 2}_{\delta(2,3) = \tuple{1,3,R}}.
\]
\end{explain}

\begin{thm}
$\Nat$ నుండి~$\Nat$కు వెళ్లే, ట్యూరింగ్-గణనీయాలు కాని ప్రమేయాలు ఉన్నాయి.
\end{thm}

\begin{proof}
ధన పూర్ణ సంఖ్యల పరిమిత క్రమాల సమితి~$(\PosInt)^*$ \tetoken{లెక్కించదగిన}{!!{enumerable}} అని
మనకు తెలుసు (\cref{sfr:siz:zigzag:prob:posint-star}). ప్రామాణిక ట్యూరింగ్
యంత్రాల వర్ణనల సమితి~$(\PosInt)^*$ యొక్క ఉపసమితి కనుక, అది కూడా లెక్కించదగినదని
తెలుస్తుంది. $\Nat$ నుండి~$\Nat$కు వెళ్లే ప్రతి ట్యూరింగ్-గణనీయ ప్రమేయాన్ని
ఏదో ఒక (నిజానికి, అనేక) ట్యూరింగ్ యంత్రం గణిస్తుంది. దాని స్థితులు,
సంకేతాలకు ధన పూర్ణ సంఖ్యలుగా పేర్లు మార్చడం ద్వారా (ప్రత్యేకించి,
$\TMendtape$కు~$1$, $\TMblank$కు~$2$, $\TMstroke$కు~$3$ అని పెట్టడం ద్వారా),
ప్రతి ట్యూరింగ్-గణనీయ ప్రమేయాన్నీ ఒక ప్రామాణిక ట్యూరింగ్ యంత్రం గణిస్తుందని
చూడవచ్చు. అంటే $\Nat$ నుండి~$\Nat$కు వెళ్లే అన్ని ట్యూరింగ్-గణనీయ ప్రమేయాల
సమితి కూడా లెక్కించదగినదే.

మరోవైపు, $\Nat$ నుండి~$\Nat$కు వెళ్లే అన్ని ప్రమేయాల సమితి \tetoken{లెక్కించదగిన}{!!{enumerable}}
కాదు (\cref{sfr:siz:red:prob:nat-nat}). అన్ని ప్రమేయాలూ ఏదో ఒక ట్యూరింగ్
యంత్రంతో గణనీయాలైతే, వాటిని గణించే ట్యూరింగ్ యంత్రాల వర్ణనలన్నిటినీ జాబితా
చేసి అన్ని ప్రమేయాల సమితినీ లెక్కించగలిగేవారం. కాబట్టి ట్యూరింగ్-గణనీయాలు
కాని కొన్ని ప్రమేయాలు ఉన్నాయి.
\end{proof}

\begin{prob}
  స్థితులను, వర్ణమాల సంకేతాలను స్పష్టంగా జాబితా చేయనవసరం లేని విధంగా
  ట్యూరింగ్ యంత్రాలను వర్ణించే ఒక పద్ధతిని మీరు ఆలోచించగలరా? ``ప్రామాణిక''
  యంత్రం అనే మీ స్వంత భావనను నిర్వచించవచ్చు; కానీ ప్రతి ట్యూరింగ్-గణనీయ
  ప్రమేయాన్నీ మీ కొత్త అర్థంలోని ఒక ``ప్రామాణిక'' యంత్రం ఎందుకు గణించగలదో
  కూడా చెప్పండి.
\end{prob}\sourcecorrection{OLTETURUND-003}{చివరి అభ్యాసంలో ఆంగ్ల మూలం ``every Turing machine can be computed by a standard machine'' అని యంత్రాన్నే గణించే వస్తువుగా పేర్కొంది. ముందు నిరూపణకు అనుగుణంగా, ప్రతి ట్యూరింగ్-గణనీయ ప్రమేయాన్ని ప్రామాణిక యంత్రం గణించగలదని చూపమనే ఉద్దేశిత ప్రశ్నను పునరుద్ధరించాం.}

\end{document}
